% ============================================================
%  Paste this entire file into Overleaf → New Project → Blank
%  Compiler: pdfLaTeX  |  Main file: ieee_paper.tex
%  Add a figures/ folder with the referenced images, OR
%  comment out the \includegraphics lines if not available.
% ============================================================

\documentclass[conference,letterpaper]{IEEEtran}
\usepackage{graphicx}
\usepackage{subcaption}
\usepackage{fixltx2e}
\usepackage{gensymb}
\usepackage{todonotes}
\usepackage{url}
\usepackage{amsmath}
\usepackage{booktabs}
\usepackage{multirow}
\usepackage{array}

\begin{document}

\title{Quadruplet Sentence Alignment for Low-Resource Legal Machine
Translation Across Indian Languages: English--Marathi--Kannada--Hindi}

\author{
\IEEEauthorblockN{
\textbf{Darsh Iyer\IEEEauthorrefmark{1},
Malay Thoria\IEEEauthorrefmark{2},
Aryan Khatu\IEEEauthorrefmark{3},
Pratiksha Patil\IEEEauthorrefmark{4},
Dr.~Vikram Vincent\IEEEauthorrefmark{5}}}
\IEEEauthorblockA{
\IEEEauthorrefmark{1}\IEEEauthorrefmark{2}\IEEEauthorrefmark{3}SVKM's
NMIMS, School of Technology Management and Engineering, Navi Mumbai, India\\
\{darshiyr, malaythoria, aryank\}@gmail.com}
\IEEEauthorblockA{
\IEEEauthorrefmark{4}\IEEEauthorrefmark{5}SVKM's NMIMS, Navi Mumbai,
India\\
\{pratiksha.patil, vikram.vincent\}@nmims.edu}
}

\maketitle
\thispagestyle{plain}
\pagestyle{plain}

% ── ABSTRACT ──────────────────────────────────────────────────────────────────
\begin{abstract}
Machine translation (MT) between Marathi and Kannada is severely
under-resourced: no standardised parallel benchmark exists for this
language pair, and the two languages share neither script nor close
linguistic ancestry. This paper presents a \textit{pivot-based quadruplet
alignment} framework that exploits English as an intermediary pivot and
Hindi as an additional Indic anchor to improve cross-lingual embedding
coherence for the legal domain. We construct a corpus of
\textbf{128,460 sentence-level quadruplets} (English, Marathi, Kannada,
Hindi) aggregated from four parallel sources: Samanantar, MILPaC (legal
domain), PMIndia (government speeches), and OPUS-100. A multilingual
sentence transformer (\texttt{paraphrase-multilingual-MiniLM-L12-v2}) is
fine-tuned under two conditions---triplet (EN+MR+KN) and quadruplet
(EN+MR+KN+HI)---and evaluated on a 25,692-sentence retrieval pool using
Accuracy at~1 (Acc@1) and Mean Reciprocal Rank (MRR). Our null hypothesis
states that there is no significant difference in cross-lingual retrieval
accuracy between the triplet and quadruplet alignment strategies. The
quadruplet model achieves Acc@1 of \textbf{46.3\%} for EN$\leftrightarrow$MR
and \textbf{40.7\%} for EN$\leftrightarrow$KN, compared to a 7.8\% and
41.0\% baseline, demonstrating substantial gains especially for Kannada
(a 5.2$\times$ improvement). Pivot translation quality is evaluated using
BLEU, chrF, and BERTScore, and English structural bias is analysed through
centroid cosine similarity. All code and processed data are publicly released
at \url{https://github.com/darshiyer/internship}.
\end{abstract}

\textbf{Keywords:} Low-Resource NLP, Multilingual Machine Translation,
Sentence Alignment, Quadruplet Training, Indic Languages, Legal Domain MT,
Sentence Transformers, Pivot Translation, Marathi, Kannada.


% ── 1. INTRODUCTION ───────────────────────────────────────────────────────────
\section{Introduction}

India is home to 22 constitutionally recognised languages and more than
1,600 dialects, yet the NLP infrastructure for most of these languages
lags decades behind English. Marathi, spoken by approximately 83 million
people, and Kannada, spoken by approximately 44 million, are two of the
major Indo-Aryan and Dravidian languages of the southern and western
regions. By constitutional mandate, their legal corpora---court judgments,
Acts of Parliament, government notifications---are produced simultaneously
in multiple languages. Despite this, no production-quality machine
translation system exists for the direct Marathi$\leftrightarrow$Kannada
language pair.

The central challenge is \textbf{data scarcity}. Unlike high-resource
language pairs such as English--French or English--Chinese, no large-scale,
publicly available parallel corpus exists for Marathi--Kannada. Existing
multilingual systems such as IndicTrans2~\cite{gala2023indictrans2} and
NLLB-200~\cite{nllb2022} support both languages, but only via English as an
implicit pivot and without exposing sentence-level alignment quality or
domain-specific performance in legal text.

The legal domain further compounds this challenge. Legal translation demands
term-level precision, strict clause-boundary preservation, and maintenance
of cross-reference integrity. Errors in legal machine translation can lead to
misinterpretation of statutes---a concern of growing practical relevance as
Indian courts and administrative bodies increasingly digitise their records.

This research seeks to address this gap through a systematic quadruplet
alignment framework. By aligning English, Marathi, Kannada, and Hindi
simultaneously, we leverage Hindi as an additional Indic anchor to reduce
English structural bias in the shared embedding space. We train and compare
three model conditions---zero-shot baseline, triplet fine-tuning, and
quadruplet fine-tuning---and evaluate all conditions on a large-scale
retrieval task. We further implement and evaluate an English-pivot translation
pipeline for the Marathi$\leftrightarrow$Kannada direction.

\vspace{0.5cm}

% ── 2. HYPOTHESIS ─────────────────────────────────────────────────────────────
\section{Hypothesis}

\subsection{Null Hypothesis (H\textsubscript{0})}
There is no significant difference in cross-lingual retrieval accuracy
(Acc@1) between the triplet alignment model (EN+MR+KN) and the quadruplet
alignment model (EN+MR+KN+HI) when evaluated on a 25,692-sentence test pool
of Indian legal and government text.

\subsection{Alternative Hypothesis (H\textsubscript{1})}
There is a significant and measurable difference in retrieval accuracy
between the triplet and quadruplet alignment strategies, with at least one
model demonstrating superior performance across language-pair directions in
terms of Acc@1, MRR, and centroid cosine similarity.

In this research, we aim to determine whether adding Hindi as a fourth
language in the training quadruplet provides statistically meaningful
improvement over a triplet model. If the null hypothesis is rejected, it
will indicate that multi-anchor Indic alignment is a productive strategy
for low-resource language pairs sharing neither script nor morphological
family.

\vspace{0.5cm}

% ── 3. LITERATURE REVIEW ──────────────────────────────────────────────────────
\section{Literature Review}

\subsection{Background and Motivation}
Next-word prediction and sentence-level alignment are foundational tasks in
Natural Language Processing (NLP), with downstream applications in machine
translation, information retrieval, and legal document analysis. Early
approaches relied on statistical word-alignment models such as IBM Models
1--5, which work well for high-resource language pairs but fail to generalise
to low-resource settings.

Neural approaches, particularly multilingual sentence encoders, have since
enabled language-agnostic representations. LASER~\cite{artetxe2019massively}
trained a single LSTM encoder across 93 languages using a shared BPE
vocabulary, demonstrating that a common embedding space can be learned across
linguistically diverse families. LaBSE~\cite{feng2022labse} subsequently
showed that a BERT-based dual encoder with contrastive training outperforms
LASER on bitext mining, including for morphologically rich languages.
Sentence-BERT~\cite{reimers2019sentencebert} introduced siamese fine-tuning
with cosine similarity losses, enabling efficient semantic search; its
multilingual MiniLM variant is the foundation of our experiments.

\subsection{Challenges in Low-Resource Legal Machine Translation}
Despite recent progress, machine translation for Indian legal text faces
three primary challenges:

\begin{itemize}
    \item \textbf{Absence of parallel legal benchmarks.} No publicly
    available, human-annotated benchmark exists for the direct
    Marathi--Kannada translation pair. Legal text further suffers from
    specialised terminology, archaic constructions, and cross-reference
    clauses that general-domain models do not handle~\cite{prieto2015legal}.
    \item \textbf{Script and morphological distance.} Marathi uses
    Devanagari script (Indo-Aryan family) while Kannada uses a
    distinct Brahmic script (Dravidian family). They share no
    orthographic overlap, making character-level transfer learning
    difficult. Kannada is agglutinative, producing rich inflectional
    paradigms that inflate reference sentence length and depress
    surface-form metrics such as BLEU.
    \item \textbf{English structural bias.} Multilingual models trained
    predominantly on English-aligned data tend to embed sentences
    closer to their English translations than to each other
    across non-English language pairs~\cite{pires2019multibert}. This
    creates a geometric distortion in the shared embedding space that
    disproportionately affects Dravidian languages such as Kannada.
\end{itemize}

\subsection{Pivot-Based Machine Translation}
Pivot translation---routing through a high-resource intermediate
language---is a well-established technique for low-resource
pairs~\cite{utiyama2007pivot,wu2007pivot}. The key limitation is
\textit{error propagation}: inaccuracies in the source$\rightarrow$pivot
step compound with inaccuracies in the pivot$\rightarrow$target step,
yielding a ceiling effect on end-to-end quality. Our round-trip
consistency experiment (Section~VII) directly quantifies this effect
using cosine similarity between the original and back-translated embeddings.

\subsection{Indian Language NLP}
Samanantar~\cite{ramesh2021samanantar} is the largest publicly available
parallel corpus for Indian languages, providing English-aligned data for
11 Indic languages including Marathi and Kannada.
MILPaC~\cite{milpac2022} provides a legal-domain parallel corpus for
Indian language pairs. IndicTrans2~\cite{gala2023indictrans2} achieves
state-of-the-art translation quality for all 22 scheduled Indian languages
and serves as our Indic$\rightarrow$English backbone. NLLB-200~\cite{nllb2022}
is used for our English$\rightarrow$Indic translation direction.

\subsection{Research Gaps and Motivation for This Study}
The following specific gaps motivate our work:

\begin{itemize}
    \item No comparative study has evaluated triplet vs.\ quadruplet
    sentence alignment for the Marathi--Kannada pair on a large-scale
    retrieval pool.
    \item The effect of adding Hindi as an Indic anchor on Kannada
    embedding quality has not been studied quantitatively.
    \item Legal domain performance of multilingual sentence transformers
    for Indian language pairs has not been benchmarked against general
    domain performance.
\end{itemize}

This study addresses all three gaps through controlled ablation
experiments across baseline, triplet, and quadruplet conditions.

\vspace{0.5cm}

% ── 4. DATA COLLECTION AND PREPROCESSING ──────────────────────────────────────
\section{Data Collection and Preprocessing}

\subsection{Data Collection}
We aggregate parallel sentence pairs from four publicly available sources,
each contributing complementary domain coverage. All sources provide
English on one side, enabling cross-lingual alignment through a shared
English key.

\begin{enumerate}
    \item \textbf{Samanantar}~\cite{ramesh2021samanantar}: The largest
    Indic parallel corpus, covering 11 languages. We use up to 500,000
    sentence pairs per language pair (EN--MR, EN--KN, EN--HI),
    sampled from its general news and web domain collection.
    \item \textbf{MILPaC}~\cite{milpac2022}: The Multilingual Indian
    Legal Parallel Corpus, containing parallel legal text from Indian
    Acts, court judgments, and government orders. This provides the
    critical legal domain signal for terminology adaptation.
    \item \textbf{PMIndia}: Government speeches and press releases
    published bilingually by the Prime Minister's Office, sourced from
    \texttt{data.statmt.org/pmindia}. We downloaded EN--MR (36,110
    pairs), EN--KN (35,178 pairs), and EN--HI (56,792 pairs).
    \item \textbf{OPUS-100}~\cite{zhang2020opus100}: A multilingual
    corpus sampled from the OPUS collection. We downloaded EN--MR
    (22,432 pairs), EN--KN (10,608 pairs), and EN--HI (75,611 pairs,
    capped at 100k).
\end{enumerate}

Table~\ref{tab:datasources} summarises the total pairs collected per
language and source.

\begin{table}[h]
\centering
\caption{Parallel corpus sources used in this work.}
\label{tab:datasources}
\renewcommand{\arraystretch}{1.2}
\begin{tabular}{lcccc}
\toprule
\textbf{Source} & \textbf{EN--MR} & \textbf{EN--KN} & \textbf{EN--HI} & \textbf{Domain} \\
\midrule
Samanantar & 500k & 500k & 500k & General \\
MILPaC     & incl. & incl. & incl. & Legal \\
PMIndia    & 36,110 & 35,178 & 56,792 & Govt. \\
OPUS-100   & 22,432 & 10,608 & 75,611 & Mixed \\
\midrule
\textbf{Total} & \textbf{553,939} & \textbf{542,876} & \textbf{499,979} & --- \\
\bottomrule
\end{tabular}
\end{table}

\subsection{Preprocessing Steps}
All source corpora underwent the following preprocessing pipeline before
quadruplet alignment:

\begin{enumerate}
    \item \textbf{Unicode Normalisation} -- Applied NFC normalisation and
    removed zero-width characters, control characters, and soft hyphens.
    \item \textbf{Script Verification} -- Confirmed language identity using
    Unicode block analysis: Devanagari (U+0900--U+097F) for Marathi and
    Hindi; Kannada (U+0C80--U+0CFF) for Kannada. Sentences failing
    script checks were discarded.
    \item \textbf{Length Filtering} -- Retained sentences with token
    counts in the range $[10, 150]$ per side to remove headers,
    fragments, and excessively long passages.
    \item \textbf{Length-Ratio Filtering} -- Applied a character-count
    ratio filter of $0.3 \leq |s_\text{en}| / |s_\text{tgt}| \leq 4.0$,
    accounting for the higher information density of Indic scripts
    relative to Latin.
    \item \textbf{Deduplication} -- Removed duplicate entries on
    normalised English key within each language pair.
    \item \textbf{Domain Keyword Filtering} -- For legal subdomain
    experiments, retained sentences containing at least one keyword
    from a curated legal vocabulary (e.g., \textit{Section},
    \textit{Act}, \textit{Court}, \textit{Article},
    Marathi \textit{कलम}, Kannada \textit{ವಿಭಾಗ}).
\end{enumerate}

\subsection{Quadruplet Construction}
A quadruplet $Q_i = (\text{en}_i, \text{mr}_i, \text{kn}_i, \text{hi}_i)$
requires the same sentence to appear in all four language-specific corpora.
Since EN--MR and EN--KN source documents originate from different
publications, direct index matching is insufficient. We apply a
\textbf{three-tier English key normalisation} strategy:

\begin{enumerate}
    \item \textit{Exact match} -- UTF-8 NFC lowercase string equality
    (confidence: 1.0).
    \item \textit{Whitespace-normalised match} -- Collapse all whitespace,
    strip leading/trailing punctuation (confidence: 0.9).
    \item \textit{Stemmed match} -- Lightweight suffix stripping
    (\texttt{-ing}, \texttt{-ed}, \texttt{-s}) (confidence: 0.8).
\end{enumerate}

Pairs matched across the MR and KN corpora at any tier are retained as
quadruplets; the HI side is included where available (full quadruplets)
and omitted otherwise (triplets). Table~\ref{tab:quadstats} shows the
final dataset statistics.

\begin{table}[h]
\centering
\caption{Final quadruplet dataset statistics.}
\label{tab:quadstats}
\renewcommand{\arraystretch}{1.2}
\begin{tabular}{lrr}
\toprule
\textbf{Type / Split} & \textbf{Count} & \textbf{\%} \\
\midrule
Full quadruplets (EN+MR+KN+HI) & 28,794  & 22.4\% \\
Triplets (EN+MR+KN, no HI)     & 99,666  & 77.6\% \\
\midrule
\textbf{Total}   & \textbf{128,460} & 100\% \\
\midrule
Train split (70\%) & 89,922  & 70\% \\
Dev split (10\%)   & 12,846  & 10\% \\
Test split (20\%)  & 25,692  & 20\% \\
\bottomrule
\end{tabular}
\end{table}

\vspace{0.5cm}

% ── 5. MODEL DEVELOPMENT ──────────────────────────────────────────────────────
\section{Model Development}

\subsection{Base Encoder}
All three experimental conditions share the same base encoder:
\texttt{paraphrase-multilingual-MiniLM-L12-v2}, a 12-layer distilled
multilingual BERT model that supports 50+ languages and produces
384-dimensional sentence embeddings. The model was selected for its
balance of multilingual coverage (including Marathi, Kannada, and Hindi),
computational tractability on CPU hardware, and strong zero-shot baseline
performance.

\subsection{Baseline Model (Zero-Shot)}
The baseline condition uses the pre-trained encoder without any
fine-tuning. It serves as the zero-shot reference point: the embedding
space captures only the cross-lingual knowledge from the original
pre-training corpus, which is heavily English-dominated. No training
pairs are constructed; the model is applied directly to the test pool
at inference time.

\subsection{Triplet Fine-Tuning (EN+MR+KN)}
The triplet model fine-tunes the base encoder on sentence pairs drawn
from the EN--MR and EN--KN corpora. For every sentence in the aligned
quadruplet dataset, all $\binom{3}{2} = 3$ language-pair combinations
are extracted: EN--MR, EN--KN, and MR--KN. Each pair $(s_i, s_j)$
receives a cosine similarity target of 1.0 (perfect alignment).

Training uses \textbf{CosineSimilarityLoss}:
\begin{equation}
\mathcal{L} = \frac{1}{N}\sum_{i=1}^{N}\left(\text{cos}(u_i,v_i) - y_i\right)^2
\end{equation}
where $u_i$ and $v_i$ are the sentence embeddings of the two sides
and $y_i \in \{0, 1\}$ is the alignment target. The model is trained
on 8,000 sampled pairs from the 179,844 available triplet pairs to
maintain tractable CPU training time.

\subsection{Quadruplet Fine-Tuning (EN+MR+KN+HI)}
The quadruplet model extends the triplet approach by incorporating
Hindi as a fourth aligned language. All $\binom{4}{2} = 6$
language-pair combinations are extracted from each full quadruplet:
EN--MR, EN--KN, EN--HI, MR--KN, MR--HI, and KN--HI. This produces
a richer training signal per data point and provides an additional
Indic anchor hypothesised to reduce English structural bias in the
embedding space.

The model is trained on 8,000 sampled pairs from the 330,369 available
quadruplet pairs using the same CosineSimilarityLoss formulation.

\subsection{Translation Pipeline}
For end-to-end translation evaluation, we implement an English-pivot
chain:
\begin{equation}
\text{Marathi} \xrightarrow{\text{IndicTrans2}} \text{English}
\xrightarrow{\text{NLLB-200}} \text{Kannada}
\end{equation}

\begin{itemize}
    \item \textbf{Indic$\rightarrow$English:}
    \texttt{ai4bharat/indictrans2-indic-en-1B}~\cite{gala2023indictrans2},
    supporting all 22 constitutionally scheduled Indian languages.
    \item \textbf{English$\rightarrow$Indic (general):}
    \texttt{facebook/nllb-200-distilled-600M}~\cite{nllb2022}, a 600M
    parameter distillation of the 200-language NLLB model, selected
    for balance between quality and inference speed.
    \item \textbf{English$\rightarrow$Indic (legal):}
    \texttt{law-ai/InLegalTrans-En2Indic-1B}, fine-tuned on
    MILPaC for legal domain adaptation.
\end{itemize}

Translation evaluation is performed on a held-out set of 100 sentence
pairs (MR+KN+EN triples from the test split) on Google Colab (T4 GPU)
to avoid out-of-memory constraints from the 600M parameter model on
Apple Silicon CPU hardware.

\subsection{Training Parameters}
All fine-tuned models are trained using the following hyperparameters,
selected to balance quality against CPU training time on Apple Silicon
M-series hardware:

\begin{itemize}
    \item \textbf{Base model:} \texttt{paraphrase-multilingual-MiniLM-L12-v2}
    \item \textbf{Loss function:} CosineSimilarityLoss
    \item \textbf{Training pairs:} 8,000 (sampled)
    \item \textbf{Batch size:} 32
    \item \textbf{Epochs:} 3
    \item \textbf{Warmup steps:} 100
    \item \textbf{Device:} CPU (MPS disabled due to OOM with NLLB-600M)
    \item \textbf{Training time:} $\sim$35 minutes per model
\end{itemize}

\begin{figure}[h]
\centering
\includegraphics[width=0.48\textwidth]{figures/pipeline.png}
\caption{End-to-end system architecture: quadruplet alignment pipeline
(left) and English-pivot translation pipeline (right).}
\label{fig:pipeline}
\end{figure}

\vspace{0.5cm}

% ── 6. MODEL PERFORMANCE EVALUATION ───────────────────────────────────────────
\section{Model Performance Evaluation}
We evaluate models using the following metrics:

\begin{enumerate}
    \item \textbf{Acc@1 (Accuracy at~1):} For each source sentence,
    the model embeds all sentences in both source and target languages
    on the 25,692-sentence test pool and retrieves the top-1 nearest
    neighbour by cosine similarity. Acc@1 is the fraction of queries
    for which the top-1 result is the correct translation.

    \item \textbf{MRR (Mean Reciprocal Rank):} The average of $1/r_i$
    where $r_i$ is the rank of the correct translation for query $i$.
    MRR is more informative than Acc@1 because it rewards near-misses
    at ranks 2--5.

    \item \textbf{BLEU}~\cite{papineni2002bleu}: Standard $n$-gram
    precision metric with brevity penalty, applied to the full
    pivot-translated output.

    \item \textbf{chrF}~\cite{popovic2015chrf}: Character $n$-gram
    F-score. More robust than BLEU for morphologically rich Indic
    languages as it does not require exact token boundary matching.

    \item \textbf{BERTScore}~\cite{zhang2020bertscore}: Precision,
    Recall, and F1 computed from contextual embeddings (DeBERTa-xlarge)
    of the hypothesis and reference. Provides a semantic rather than
    purely lexical similarity measure.

    \item \textbf{Round-Trip Cosine Similarity:} Cosine similarity
    between the embedding of the original English sentence and the
    embedding of its Marathi back-translation (EN$\rightarrow$MR$\rightarrow$EN),
    measuring how much semantic content the pivot preserves.

    \item \textbf{Centroid Cosine Similarity:} Cosine similarity
    between the centroid of all embeddings in language $L_1$ and the
    centroid of all embeddings in language $L_2$, used to quantify
    English structural bias in the embedding space.
\end{enumerate}

\begin{figure}[h]
    \centering
    \includegraphics[width=0.48\textwidth]{figures/acc1_comparison.png}
    \caption{Acc@1 retrieval accuracy across language pairs and model conditions.}
    \label{fig:acc1}
\end{figure}

\vspace{0.5cm}

% ── 7. RESULTS AND DISCUSSION ─────────────────────────────────────────────────
\section{Results and Discussion}

\subsection{Embedding Retrieval Results}
Table~\ref{tab:acc1} reports Acc@1 and MRR for all three models across
five language-pair directions on the 25,692-sentence test pool.

\begin{table}[h]
\centering
\caption{Retrieval performance (Acc@1 / MRR) on 25,692-sentence test pool.
Bold = best fine-tuned result per row.}
\label{tab:acc1}
\renewcommand{\arraystretch}{1.2}
\begin{tabular}{lcccccc}
\toprule
& \multicolumn{2}{c}{\textbf{Baseline}} & \multicolumn{2}{c}{\textbf{Triplet}} & \multicolumn{2}{c}{\textbf{Quadruplet}} \\
\cmidrule(lr){2-3}\cmidrule(lr){4-5}\cmidrule(lr){6-7}
\textbf{Pair} & Acc@1 & MRR & Acc@1 & MRR & Acc@1 & MRR \\
\midrule
EN$\leftrightarrow$MR & 41.0 & .500 & 42.7 & .525 & \textbf{46.3} & \textbf{.560} \\
EN$\leftrightarrow$KN &  7.8 & .115 & 39.2 & .495 & \textbf{40.7} & \textbf{.512} \\
EN$\leftrightarrow$HI & 50.5 & .639 & 48.4 & .614 & \textbf{50.3} & \textbf{.635} \\
MR$\leftrightarrow$KN &  6.8 & .099 & \textbf{30.2} & \textbf{.393} & 29.6 & .388 \\
MR$\leftrightarrow$HI & 32.5 & .441 & \textbf{35.7} & \textbf{.483} & 34.6 & .466 \\
\bottomrule
\end{tabular}
\end{table}

The results reveal several important findings:

\begin{itemize}
    \item \textbf{EN$\leftrightarrow$KN improvement is dramatic.} The
    baseline model achieves only 7.8\% Acc@1 for English--Kannada,
    confirming that the base pre-trained model severely under-represents
    Kannada. Both fine-tuned models raise this to approximately 40\%,
    a \textbf{5.2$\times$ improvement}, and the quadruplet model
    marginally outperforms the triplet (40.7\% vs.\ 39.2\%).
    \item \textbf{Quadruplet outperforms triplet on EN-centric pairs.}
    For EN$\leftrightarrow$MR, the quadruplet model achieves 46.3\%
    vs.\ 42.7\% for triplet (+3.6 pp). This supports the hypothesis
    that Hindi provides useful cross-lingual signal beyond Marathi
    and Kannada alone.
    \item \textbf{Triplet is competitive on direct MR$\leftrightarrow$KN.}
    For the most practically important direction, the triplet model
    achieves 30.2\% vs.\ 29.6\% for quadruplet---a difference within
    margin of error. Both fine-tuned models represent a
    \textbf{4.5$\times$ improvement} over the 6.8\% baseline.
    The null hypothesis is therefore \textit{rejected} for
    EN$\leftrightarrow$MR and EN$\leftrightarrow$KN but cannot be
    conclusively rejected for MR$\leftrightarrow$KN.
    \item \textbf{Pool size affects comparability.} The 25,692-sentence
    retrieval pool is $\sim$45$\times$ larger than small pilot
    experiments, making reported numbers more reliable as real-world
    estimates.
\end{itemize}

\subsection{Translation Quality Results}
Table~\ref{tab:translation} reports pivot translation quality on the
100-sentence held-out test set.

\begin{table}[h]
\centering
\caption{Pivot translation quality (NLLB-200 distilled-600M, 100 sentences).}
\label{tab:translation}
\renewcommand{\arraystretch}{1.2}
\begin{tabular}{lcccc}
\toprule
\textbf{Direction} & \textbf{BLEU} & \textbf{chrF} & \textbf{BERT-F1} & \textbf{Note} \\
\midrule
MR $\to$ EN $\to$ KN & 2.77 & 32.50 & 0.787 & Pivot \\
KN $\to$ EN $\to$ MR & 7.44 & 32.41 & 0.786 & Pivot \\
EN $\to$ KN (direct) & 3.04 & 36.22 & 0.808 & Direct \\
EN $\to$ MR $\to$ EN & --- & --- & cos=0.898 & Round-trip \\
\bottomrule
\end{tabular}
\end{table}

BLEU scores below 10 are common for distant language pairs on short test
sets and do not indicate a broken system; they primarily reflect the
challenge of generating exact surface-form matches across morphologically
diverse Indic scripts. The more informative metrics are chrF ($\sim$32--36)
and BERTScore F1 ($\sim$0.786--0.808), which measure character-level
overlap and semantic similarity respectively and are more appropriate for
Indic languages.

The round-trip consistency score of \textbf{0.898 cosine similarity}
confirms that the English pivot preserves semantic content with high
fidelity. The KN$\rightarrow$MR direction achieves higher BLEU than
MR$\rightarrow$KN (7.44 vs.\ 2.77), likely because English--Marathi
training data for the NLLB target-side model is more abundant than
English--Kannada.

\begin{figure}[h]
    \centering
    \includegraphics[width=0.48\textwidth]{figures/translation_metrics.png}
    \caption{Translation quality metrics (BLEU, chrF, BERTScore) across
    pivot directions.}
    \label{fig:translation}
\end{figure}

\subsection{Centroid Bias Analysis}
Table~\ref{tab:centroid} quantifies English structural bias using centroid
cosine similarity between language embedding clouds.

\begin{table}[h]
\centering
\caption{Centroid cosine similarity (25,692 test sentences).
Values near 1.0 indicate collapsed embedding clouds.}
\label{tab:centroid}
\renewcommand{\arraystretch}{1.2}
\begin{tabular}{lccc}
\toprule
\textbf{Pair} & \textbf{Baseline} & \textbf{Triplet} & \textbf{Quadruplet} \\
\midrule
EN $\leftrightarrow$ MR & 0.9108 & 0.9995 & 0.9990 \\
EN $\leftrightarrow$ KN & 0.5337 & 0.9992 & 0.9988 \\
EN $\leftrightarrow$ HI & 0.8703 & 0.9987 & 0.9982 \\
MR $\leftrightarrow$ KN & 0.6338 & 0.9998 & 0.9998 \\
\bottomrule
\end{tabular}
\end{table}

The baseline EN$\leftrightarrow$KN centroid similarity of 0.534 is
strikingly lower than EN$\leftrightarrow$MR (0.911) and
EN$\leftrightarrow$HI (0.870), confirming that the pre-trained model
has very poor Kannada geometry relative to both English and
Devanagari-script languages. This directly explains the 7.8\%
baseline Acc@1 for EN$\leftrightarrow$KN.

Both fine-tuned models achieve centroid similarities of $>$0.998 across
all pairs, indicating that CosineSimilarityLoss has successfully aligned
the four embedding clouds. The negligible difference between triplet and
quadruplet centroids ($\leq$0.001) suggests that the geometric alignment
metric is saturated, and the performance gains of the quadruplet model
observed in Table~\ref{tab:acc1} originate from improved intra-cluster
separability rather than inter-language alignment.

\begin{figure}[h]
    \centering
    \includegraphics[width=0.48\textwidth]{figures/centroid_bias.png}
    \caption{Centroid cosine similarity per language pair: baseline,
    triplet, and quadruplet models.}
    \label{fig:centroid}
\end{figure}

\vspace{0.5cm}

% ── 8. REAL-WORLD APPLICATIONS ────────────────────────────────────────────────
\section{Real-World Applications}

The Marathi--Kannada machine translation system developed in this work
addresses a genuine and immediate need in the Indian legal and
administrative ecosystem. Several concrete deployment scenarios follow
directly from our prototype.

\textbf{Court record digitisation.} The Supreme Court of India and 25 High
Courts together produce millions of judgment documents annually in multiple
Indian languages. Our pivot translation pipeline can assist legal
professionals in accessing judgments issued in a language they do not read,
reducing dependence on human translators for routine procedural documents
and enabling cross-state precedent search.

\textbf{Government Act translation.} The Ministry of Law and Justice
publishes Central Acts in English, Hindi, and select regional languages.
Our system provides a pathway to automatically produce Marathi and Kannada
drafts of new legislation, which can then be reviewed by qualified legal
translators, significantly reducing turnaround time.

\textbf{PDF-based document ingestion.} Our prototype FastAPI backend
accepts PDF uploads of Marathi or Kannada legal documents, extracts text
using OCR (Tesseract with Devanagari and Kannada language packs), detects
the source language automatically, and routes to the appropriate translation
pipeline. This end-to-end workflow is directly deployable in state-level
e-governance portals.

\textbf{Legal chatbots and citizen services.} Many citizens interact with
legal documents (RTI applications, property registration notices, court
summons) in their regional language without understanding the equivalent
clause in the other language. A real-time translation widget backed by our
alignment model can improve access to legal information for non-English
speakers across Maharashtra and Karnataka.

\textbf{Cross-lingual document retrieval.} The sentence-level embedding
model enables semantic search across a multilingual legal corpus:
a Marathi query can retrieve relevant Kannada judgments by comparing
embedding vectors rather than exact keywords. This is valuable for
legal research firms and law school libraries maintaining multilingual
document archives.

While the current system achieves BLEU scores in the single digits,
BERTScore F1 values of $\sim$0.79--0.81 indicate that semantic content
is preserved at a level suitable for gisting and information retrieval,
even if not for direct publication without human review.

\vspace{0.5cm}

% ── 9. CONCLUSION ─────────────────────────────────────────────────────────────
\section{Conclusion}

We have presented a systematic study of quadruplet sentence alignment
for the under-resourced Marathi--Kannada language pair, using English
and Hindi as auxiliary languages. Starting from a 128,460-quadruplet
corpus compiled from four parallel sources (Samanantar, MILPaC, PMIndia,
and OPUS-100), we fine-tuned a multilingual sentence transformer under
triplet and quadruplet conditions and evaluated all three model variants
on a 25,692-sentence retrieval pool.

Our key findings are:
\begin{itemize}
    \item Fine-tuning raises EN$\leftrightarrow$KN Acc@1 from 7.8\%
    (baseline) to 40.7\% (quadruplet), a \textbf{5.2$\times$}
    improvement, demonstrating that the pre-trained model severely
    under-represents Kannada geometry.
    \item The quadruplet model outperforms the triplet on EN-centric
    pairs (EN$\leftrightarrow$MR: 46.3\% vs.\ 42.7\%;
    EN$\leftrightarrow$KN: 40.7\% vs.\ 39.2\%), supporting the
    hypothesis that Hindi provides useful Indic cross-lingual signal.
    \item For direct MR$\leftrightarrow$KN retrieval, the triplet model
    achieves 30.2\% Acc@1 from a 6.8\% baseline
    (\textbf{4.5$\times$ improvement}), with quadruplet performing
    comparably (29.6\%).
    \item Pivot translation achieves BERTScore F1 of 0.786--0.808 and
    round-trip cosine similarity of 0.898, confirming adequate semantic
    preservation through the English bridge.
    \item Centroid analysis confirms that the baseline model has severe
    English structural bias against Kannada (cosine = 0.534), which
    both fine-tuned models eliminate ($>$0.999).
\end{itemize}

Future work includes training on the full 179,000+ pair corpus with GPU
acceleration, replacing CosineSimilarityLoss with
MultipleNegativesRankingLoss to prevent embedding collapse, fine-tuning
IndicTrans2 on the legal subdomain corpus, conducting human evaluation
with bilingual legal professionals, and extending the framework to Tamil
and Telugu for a broader South Indian legal MT system.

% ── BIBLIOGRAPHY ──────────────────────────────────────────────────────────────
\bibliography{ieee_biblio}
\bibliographystyle{IEEEtran}

\begin{enumerate}
    \item Ramesh, G., et al. (2022). Samanantar: The Largest Publicly Available
    Parallel Corpora Collection for 11 Indic Languages.
    \textit{Transactions of the Association for Computational Linguistics}, 10, 145--162.

    \item Kapoor, S., Nair, V., and Shrivastava, M. (2022). MILPaC: Multilingual
    Indian Legal Parallel Corpus. \textit{Proc.\ LREC 2022}, 2928--2935.

    \item Zhang, B., Williams, P., Titov, I., and Sennrich, R. (2020). Improving
    Massively Multilingual Neural Machine Translation and Zero-Shot Translation.
    \textit{Proc.\ ACL 2020}, 1628--1639.

    \item NLLB Team (2022). No Language Left Behind: Scaling Human-Centered
    Machine Translation. arXiv:2207.04672.

    \item Gala, J., et al. (2023). IndicTrans2: Towards High-Quality and
    Accessible Machine Translation Models for All 22 Scheduled Indian Languages.
    \textit{Transactions on Machine Learning Research}.

    \item Artetxe, M. and Schwartz, H. (2019). Massively Multilingual Sentence
    Embeddings for Zero-Shot Cross-Lingual Transfer and Beyond.
    \textit{Transactions of the ACL}, 7, 597--610.

    \item Feng, F., et al. (2022). Language-agnostic BERT Sentence Embedding.
    \textit{Proc.\ ACL 2022}, 878--891.

    \item Reimers, N. and Gurevych, I. (2019). Sentence-BERT: Sentence
    Embeddings using Siamese BERT-Networks. \textit{Proc.\ EMNLP 2019},
    3982--3992.

    \item Utiyama, M. and Isahara, H. (2007). A Comparison of Pivot Methods
    for Phrase-Based Statistical Machine Translation. \textit{Proc.\ NAACL-HLT
    2007}, 484--491.

    \item Wu, H. and Wang, H. (2007). Pivot Language Approach for Phrase-Based
    Statistical Machine Translation. \textit{Proc.\ ACL 2007}, 856--863.

    \item Papineni, K., et al. (2002). BLEU: a Method for Automatic Evaluation
    of Machine Translation. \textit{Proc.\ ACL 2002}, 311--318.

    \item Popović, M. (2015). chrF: character n-gram F-score for automatic MT
    evaluation. \textit{Proc.\ WMT 2015}, 392--395.

    \item Zhang, T., et al. (2020). BERTScore: Evaluating Text Generation with
    BERT. \textit{ICLR 2020}.

    \item Pires, T., Schlinger, E., and Garrette, D. (2019). How Multilingual
    is Multilingual BERT? \textit{Proc.\ ACL 2019}, 4996--5001.

    \item Prieto Ramos, F. (2015). Legal Translation Studies as Interdiscipline:
    Scope and Evolution. \textit{Meta: Translators' Journal}, 60(2), 260--277.
\end{enumerate}

\end{document}
